\documentclass[11pt]{article}
\usepackage[a4paper,margin=1in]{geometry}
\usepackage{amsmath,amssymb,amsthm}
\usepackage{mathtools}
\usepackage[T1]{fontenc}
\usepackage{lmodern}
\usepackage{microtype}

\newtheorem{theorem}{Theorem}
\newtheorem{lemma}{Lemma}

\begin{document}

\title{Maximizing a Rational Objective Under Inequality Constraints}
\author{}
\date{}
\maketitle

\section*{Problem statement: no job has zero processing time}

Let
\[
Q := p_1^2+p_3^2+p_2^2+p_4^2,
\]
and let \(\alpha\) be a parameter with \(0 \le \alpha \le \tfrac12\).
We want to maximize
\[
\lambda(p_1,p_3,p_2,p_4)(\alpha)
:=
\frac{{(p_1+p_2)}^2+{(p_4+p_3)}^2+(1-2(1-\alpha))Q}{{(p_2+p_3)}^2+{(p_4+p_1)}^2+(1-2(1-\alpha))Q}
\]
over \(p_1,p_3,p_2,p_4 > 0\), subject to
\begin{align*}
&\text{(A)} \quad p_1 \le p_3+p_4,\\
&\text{(B)} \quad p_3 \le p_1+p_2,\\
&\text{(C)} \quad p_3 \le p_1+p_4,\\
&\text{(D)} \quad p_2 \le p_1+p_4,\\
&\text{(E)} \quad p_1 \le p_2+p_3,\\
&\text{(F)} \quad p_4 \le p_2+p_3,\\
\end{align*}

\section*{Statement of the result}

\begin{theorem}
If \(0<\alpha\le \tfrac12\), then the supremum over the strictly positive feasible region is
\[
\boxed{
1+\frac{1}{1+6\alpha}.
}
\]
The supremum is not attained when \(p_1,p_2,p_3,p_4>0\). It is approached, for example, by
\[
p_1=p_4=t,\qquad p_2=2t,\qquad p_3\downarrow 0,
\qquad t>0.
\]
If \(\alpha=0\), the supremum is \(2\), again approached at the same boundary.
\end{theorem}

\section*{Step 1: Simplifying the objective}

Since
\[
1-2(1-\alpha)=2\alpha-1,
\]
the numerator becomes
\begin{align*}
&{(p_1+p_2)}^2+{(p_4+p_3)}^2+(2\alpha-1)Q\\
&=2p_1p_2+2p_4p_3+2\alpha Q.
\end{align*}
Similarly,
\[
{(p_2+p_3)}^2+{(p_4+p_1)}^2+(2\alpha-1)Q
=2p_2p_3+2p_4p_1+2\alpha Q.
\]
Therefore
\[
\lambda
=
\frac{\alpha Q+p_1p_2+p_4p_3}{\alpha Q+p_1p_4+p_2p_3}.
\]
Subtracting \(1\), we get
\begin{align*}
\lambda-1
&=
\frac{p_1p_2+p_4p_3-p_1p_4-p_2p_3}
{\alpha Q+p_1p_4+p_2p_3}\\
&=
\frac{(p_1-p_3)(p_2-p_4)}
{\alpha Q+p_1p_4+p_2p_3}.
\end{align*}


\section*{Step 2: The relevant sign case}

To get a value above \(1\), we need
\[
(p_1-p_3)(p_2-p_4)>0.
\]
We treat the case
\[
p_1\ge p_3,\qquad p_2\ge p_4.
\]
The other positive-product case of \((p_1-p_3) < 0 \quad \text{and} \quad(p_2-p_4) < 0\) is symmetric and gives
the same supremum as every instance with this property may be trainsformed into case 1 via \(T(1,2,3,4) = (3,4,1,2)\).
Under which, both the objective aswell as the constraints are invariant, and the sign case is switched.
Set
\[
x:=p_1-p_3>0,\qquad y:=p_2-p_4>0.
\]
Thus
\[
p_1=p_3+x,\qquad p_2=p_4+y.
\]

\section*{Step 3: Reducing the constraints}

We now rewrite the feasibility constraints in terms of
\[
p_1=p_3+x,\qquad p_2=p_4+y,
\qquad x,y>0.
\]
The point of this step is to separate the variables \(x,y\), which determine the numerator of \(\lambda-1\), from the remaining variables \(p_3,p_4\), which will later be minimized in the denominator.

First, constraint (A) becomes
\[
p_3+x\le p_3+p_4,
\]
and hence
\[
x\le p_4.
\]
Constraint (D) becomes
\[
p_4+y\le p_3+x+p_4,
\]
and hence
\[
y\le p_3+x.
\]
These two inequalities are therefore genuine restrictions:
\[
\text{(A)}\quad x\le p_4,
\qquad
\text{(D)}\quad y\le p_3+x,
\]

It remains to check that no other constraint adds a stronger condition. Constraint (B) gives
\[
p_3\le p_3+x+p_4+y,
\]
which is automatic. Constraint (C) gives
\[
p_3\le p_3+x+p_4,
\]
which is automatic as well. Constraint (E) gives
\[
p_3+x\le p_4+y+p_3,
\]
or \(x\le p_4+y\), which is weaker than \(x\le p_4\). Constraint (F) gives
\[
p_4\le p_4+y+p_3,
\]
which is automatic.

Therefore, in this sign case the constraints relevant for the upper-bound argument are precisely
\[
x\le p_4,\qquad y\le p_3+x.
\]

\section*{Step 4: Reformulating the denominator}

For fixed \(x,y>0\), the gain term is \(xy\). The denominator is
\begin{align*}
D(p_3,p_4;x,y)
&:=\alpha Q+p_1p_4+p_2p_3\\
&=\alpha\bigl({(p_3+x)}^2+p_3^2+{(p_4+y)}^2+p_4^2\bigr)
 +(p_3+x)p_4+(p_4+y)p_3.
\end{align*}
It is strictly increasing in both \(p_3\) and \(p_4\) on the nonnegative region, since
\begin{align*}
\frac{\partial D}{\partial p_3}
&=\alpha(4p_3+2x)+2p_4+y>0,\\
\frac{\partial D}{\partial p_4}
&=\alpha(4p_4+2y)+2p_3+x>0.
\end{align*}
Hence, for fixed \(x,y\), the largest value of \(\lambda\) is obtained by minimizing \(p_3,p_4\) subject to the reduced constraints.

\section*{Step 5: Boundary cases}

Since the denominator is increasing in both \(p_3\) and \(p_4\), the smallest feasible choice is obtained by taking \(p_4=x\) and taking \(p_3\) as small as constraint (D) allows. Thus
\[
(p_3,p_4)=
\begin{cases}
(0,x), & y\le x,\\
(y-x,x), & y\ge x.
\end{cases}
\]
In the strictly positive problem, \(p_3=0\) is only a boundary value. It is nevertheless enough for the supremum, since one may replace \(p_3=0\) by \(p_3=\varepsilon>0\) and let \(\varepsilon\downarrow 0\).

If \(y\le x\), put \(0<u=y/x\le 1\). Then
\[
\lambda-1
\le
\frac{u}{3\alpha+1+2\alpha u+\alpha u^2}.
\]
This expression is increasing on \(0<u\le 1\), since its derivative is
\[
\frac{\alpha +1+ \alpha u^2}{{(3\alpha+1+2\alpha u+\alpha u^2)}^2}.
\]
Thus its maximum on this interval is attained at \(u=1\), giving
\[
\lambda-1\le \frac{1}{1+6\alpha}.
\]
If \(y\ge x\), put \(0<u=x/y\le 1\). Then
\[
\lambda-1
\le
\frac{u}{(3\alpha-1)u^2+u+3\alpha+1}.
\]
This expression is increasing on \(0<u\le 1\), because its derivative has numerator
\[
3\alpha+1-(3\alpha-1)u^2>0
\]
for \(0\le\alpha\le\tfrac12\). Hence its maximum is again attained at \(u=1\), giving
\[
\lambda-1\le \frac{1}{1+6\alpha}.
\]

\section*{Step 6: Supremum}

Combining the boundary cases,
\[
\lambda-1\le \frac{1}{1+6\alpha},
\]
so
\[
\lambda\le 1+\frac{1}{1+6\alpha}.
\]
The bound is approached by taking
\[
p_1=p_4=1,\qquad p_2=2,\qquad p_3=\varepsilon,
\]
and letting \(\varepsilon\downarrow 0\).

\end{document}
